% Capítulo: Conclusion
% Codificación: UTF-8

\chapter{Conclusion and Future Work}
\label{chap:conclusion}

The objective of this thesis was to develop an executable specification
of PISA in ASL, one that reads as documentation of the instruction set
and runs as its reference, covering a subset that models execution at
the sub-group level, with its own reconvergence model where PISA leaves
divergent control flow open, together with the tools to run kernels on
it and a test suite to check it. All of it was achieved. The
specification covers the subset of Section~\ref{sec:pisa-subset},
compiles to a C library, runs kernels compiled from OpenCL over any
number of work-groups and sub-groups, and passes every test of the
suite built to check it, against its references.

The first contribution is the specification itself. It defines the
state of a sub-group, the semantics of each instruction of the subset
as a function, and an execution model that applies one instruction per
step across the active lanes, in two passes, so that every read of a
step happens before any of its writes. Where PISA does not say how the
work-items of a sub-group execute, the specification adopts lockstep
execution, and where PISA does not say how divergent lanes reconverge,
it defines a stackless rule: at the start of every step the global
program counter is the lowest counter among the lanes that have not
terminated, and the lanes at that counter are the active ones. The rule
takes a few lines of ASL, and the active mask follows from the per-lane
program counters alone. The same source reads as a description of what
each instruction does and compiles to a library that runs it. To our
knowledge, it is the only executable specification of PISA.

The second contribution is the tooling that makes the specification
usable. PISA has no binary encoding, so a parser takes the place of a
decoder and translates a PISA kernel into the ASL function that sets up
its variables and its instruction buffer. A host library exposes the
compiled specification through the OpenCL interface, so that an
ordinary OpenCL host program allocates buffers, sets arguments,
launches a kernel over an NDRange and reads the results back. Together
they let a kernel compiled from OpenCL run on the specification without
change.

The third contribution is the test suite and the evaluation against it.
The 26 smoke tests, with around 300 cases, check every variant of every
instruction of the subset against the PISA document, and the 23 kernels
compiled from OpenCL check that the instructions compose, that the
sub-group diverges and reconverges as the model specifies, and that
data goes from the host to the kernel and back. Every case and every
kernel matches its expected output. That includes the two kernels
whose outputs depend on which lanes are active at a shuffle or a
reduction. The suite was built alongside the specification and ran
after every change. It caught the bugs of the first versions, mainly in
the handling of divergence, and it found two errors in the examples of
the PISA document, which were reported to its maintainers.

The evaluation also showed what the specification is for beyond this
thesis. The smoke tests are regression tests. If the ASL lives alongside the
PISA document, a change to an instruction is made in both, and the
tests say at once whether anything else broke. They also run the
document's own examples, which is how the two errors in them were identified.
The kernel tests also serve to check a compiler: a kernel's expected
output comes from its OpenCL source, and its actual output from the
compiled PISA run on the specification, so with the specification as
the reference, a mismatch points at the compiler. And any other tool that runs PISA kernels can be compared
with the specification by running the same kernels on the same inputs,
with no expected output needed at all.

\section{Limitations}

The specification covers a subset of PISA. It has no floating point,
no function calls, and no barriers, fences or atomics. Without floating
point it supports no floating-point workload. Without calls it cannot
run a kernel whose functions a compiler does not inline. Without
synchronization, the sub-groups of a launch run one after another, so
communication between sub-groups through memory sees the result of
that fixed order, and a workload in which sub-groups exchange data can
be racy.

PISA does not define reconvergence, so the model is a choice of the
specification, and other valid choices exist. The suite cannot tell
them apart: in the two kernels whose outputs depend on reconvergence,
the paths of every branch rejoin after both of them in instruction
order, as compiler-emitted code usually does, and on such code every common
scheme activates the same lanes at the same instructions. The rule also
orders lanes by program counter alone, which stops working once calls
are added, since a lane inside a called function can be at a lower
address than a lane that has not made the call.

Some behaviors the specification defines are not exercised by the
suite, and by design. The specification expects well-formed kernels, as
produced from real workloads by a compiler, and such kernels do not
contain unstructured branches, do not end other than in a return, and
do not store to one address from several lanes in one step. The
specification gives each of these cases a behavior, since it must do
something: it writes results back in lane order, it reconverges lanes
later or elsewhere when a join precedes a path, and it assumes that the
last instruction of a kernel is a return. It also places no bound on
the number of steps, so a kernel in which a lane never terminates does
not complete. The evaluation makes no claim about any of these cases.

The evaluation is a set of tests, so it shows that the specification
is correct on the kernels of the suite, not on every program PISA
admits. And since all the kernels were compiled with the same
pre-release build of the LLVM PISA compiler, the suite covers the code
patterns that build produces and no others.

Finally, the parser assumes a well-formed kernel, with every variable
declared and every label defined. The host
library accepts work-groups whose size is a multiple of 32 and one
kernel launch per host program.

\section{Future Work}

The first direction is extending the subset. Most of the instructions
PISA adds beyond the subset, the floating point ones first, are new
semantic functions in the existing framework and need no change to the
execution model.

The second is extending the execution model. Calls need a call stack
per lane and a reconvergence rule that accounts for call depth, since
program counter order alone no longer says which lanes are behind once
lanes are inside different calls. Collange's ordering by stack depth
before program counter~\cite{collange-stackless} is one alternative.
Barriers, fences and atomics need several sub-groups in flight at once,
with more state, a rule for scheduling between them, and a mechanism
for lanes to wait on one another. With that in place, the three have
an order to define, at the sub-group and at the work-group level.

The third is extending the test suite. Well-known test suites for GPU
algorithms, such as reductions, scans, sorting and matrix operations,
can be added once the subset supports what they need. Hand-written
PISA kernels can cover the layouts compiled code does not produce,
where the reconvergence schemes differ. And once the LLVM PISA compiler
is released upstream, the kernels of Appendix~\ref{app:kernels} can be
recompiled with it and the suite run again, to confirm that the results
hold on the released compiler.

The fourth is using the existing ASL toolchain. ASL translates into Sail, and
from Sail into emulators and into definitions for theorem
provers~\cite{isa-semantics}, where properties of the execution model
can be proved, such as the guarantee of Section~\ref{sec:asl-divergence}
that the active lanes at every step are exactly the non-terminated
lanes at the lowest counter.

The remaining directions are engineering. The parser can be completed
to accept all of PISA, so that a kernel outside the subset is reported
by the specification rather than rejected by the parser. One
observation applies here. The parser exists only because PISA has no
binary encoding. If PISA ever had one, the parser would go away: ASL
decodes instructions from their bit patterns, as the other
specifications written in it do, so the decoder would be part of the
specification itself, and supporting a new instruction would mean
adding its encoding and its semantics in one place. The host library
can accept several launches and several kernels per program, and load
kernels at run time instead of by registration. And the parts of the
specification that mirror the PISA document, the instruction tables
and the operand formats that the parser and the instruction record now
encode by hand, can be generated from the sources of the document,
which would keep the specification in step with PISA as it evolves.